Sentiment classification of social-media text is one of the most studied tasks in applied natural language processing, and comparisons between classical machine-learning classifiers and neural networks on Twitter data are abundant. Such comparisons are more sensitive to evaluation practice than is usually acknowledged: a vocabulary fitted on the test data, a fixed decision threshold, a model selected on the split that is then reported, or accuracy quoted on a corpus in which one class holds 93\,percent of the documents can each move a result by more than the difference between the methods being compared. The author's 2022 undergraduate thesis compared naive Bayes, logistic regression, a support vector machine, a random forest and gradient-boosted trees on two Twitter corpora and exhibited every one of these weaknesses.

This thesis re-implements that study as a benchmark under one evaluation protocol and extends it in three directions. Eight model families (multinomial naive Bayes, logistic regression, a calibrated linear support vector machine, random forest, XGBoost, a soft-voting ensemble, a convolutional neural network and a bidirectional LSTM) are evaluated across five input representations (count and TF-IDF $n$-grams, mean-pooled Word2Vec, Doc2Vec and token sequences) on three corpora: racist and sexist tweet detection (31,962 tweets, 7\,percent positive), COVID-19 vaccine tweets with lexicon-derived labels (10,353 tweets, three classes) and the IMDB movie-review corpus (50,000 reviews, balanced). Every one of the 56 cells of this grid follows the same procedure: representations are fitted on cross-validation training folds only, out-of-fold probabilities determine the decision threshold and the selected model, and the held-out test split is scored exactly once.

The results separate two things that fixed-threshold comparisons conflate. On the hate-speech corpus all eight families rank the test tweets almost equally well (ROC-AUC 0.939 to 0.952), while F1 on the positive class spans 0.608 to 0.684 because of calibration and threshold transfer; tuning the threshold on out-of-fold predictions is worth up to 0.09 F1 for the classical models and 0.14 for the networks, and the class-weighted networks, although best by cross-validated F1, are the least calibrated and transfer their threshold worst. On the vaccine corpus the CNN reaches macro F1 0.812 against 0.775 for the best classical model. On IMDB, linear models on TF-IDF (F1 0.884 to 0.888, accuracy 0.882 to 0.885) are the ceiling of the grid, matching published results for models without pretrained language representations, and networks trained from scratch do not reach them. Learned document embeddings are the weakest representation on every corpus. Against the 2022 study, the ranking of the classical models survives and their absolute numbers do not.

The benchmark is delivered as a maintained system: configuration-defined experiments, experiment tracking, exported pipelines with model cards and monitoring baselines, a prediction service over a relational store, population-stability drift monitoring, continuous integration and infrastructure definitions for cloud deployment. Every number and figure in this document is regenerated from saved probabilities with one command. The final chapter sets out the research programme that follows from these results in 2026: pretrained and large language models as upper references and as annotators, cross-validation bagging and conformal prediction for calibrated decisions, human annotation to replace lexicon labels, fairness auditing of hate-speech classifiers, and drift-triggered retraining.

\medskip
\noindent\textbf{Keywords:} sentiment analysis, hate speech detection, text classification, evaluation methodology, data leakage, threshold selection, calibration, reproducibility, MLOps.
